\section{Effective continuation, connection formulas, and global claims}\label{sec:continuation}

\subsection{Taylor stepping along a coded path}

A qualitative statement that a function ``can be analytically continued along
any path avoiding the singularities'' must be replaced by a finite path and
quantitative separation data.

\begin{definition}[Continuation certificate]\label{def:continuation-cert}
A continuation certificate from \(z_0\) to \(z_m\) consists of rational or
complex-rational centers
\(z_0,z_1,\ldots,z_m\), positive rational radii \(R_j\), and margins
\(d_j>0\) such that
\begin{align*}
 \norm{z_{j+1}-z_j}&\le R_j/2,\\
 \operatorname{dist}(\overline{D(z_j,R_j)},\{0,1\})&\ge d_j,
\end{align*}
plus initial names for \(y(z_0),y'(z_0)\).  If powers occur, the certificate
also contains compatible logarithm branches on the disks and overlap
constants.
\end{definition}

On a disk avoiding \(0,1\), rewrite the equation as
\[
 y''=P(z)y'+Q(z)y,
 \quad
 P(z)=-\frac{c-(a+b+1)z}{z(1-z)},
 \quad
 Q(z)=\frac{ab}{z(1-z)}.
\]
Expand \(P,Q\) at the center.  Their coefficients have geometric bounds from
the separation margin.  The Taylor coefficients of \(y\) then satisfy a
finite convolution recurrence determined by \(y(z_j),y'(z_j)\).

\begin{theorem}[Primitive recursive path continuation]\label{thm:path-continuation}
Given \cref{def:continuation-cert}, a primitive recursive algorithm computes
names for \(y(z_j),y'(z_j)\) at every center and hence a name at the endpoint.
The result is independent, up to effective equality, of refinements of the
same certified disk chain.  Applied to the regular series at zero after an
initial center inside \(|z|<1\), it gives an effective branch of
\(\F a b c z\) along the path.
\end{theorem}

\begin{proof}
For one step, majorize the coefficient recurrence by a scalar series
\(M/(1-w/R)\) whose rational constants are computed from
\(d_j,R_j,\norm a,\norm b,\norm c\) and the current value bounds.  Induction
shows that the Taylor coefficient of order \(n\) is at most \(MR^{-n}\).
At a point with displacement at most \(R/2\), the tail after \(N\) terms is
at most \(2M2^{-N}\), and the differentiated tail has a similarly explicit
bound.  Choose \(N\) by bounded search, evaluate the finite Taylor polynomial,
and propagate the resulting intervals to the next center.  Repeat for the
finite list of disks.

If a disk chain is refined, both constructions solve the same local initial
value problem on each overlap.  The coefficient recurrence proves uniqueness,
and the two tail moduli transfer the coefficient equality to equality of
names.  This is an induction over the finite refinement, not a compactness
argument over all paths.
\end{proof}

\begin{corollary}[Branch cut convention]
A conventional principal branch on
\(\C\setminus[1,\infty)\) is represented by a family of rational polygonal
paths from a base point, each carrying a positive distance from the cut and
from \(0,1\).  Equality of two evaluations is proved when their paths are
connected by an explicitly coded homotopy grid whose cells stay a positive
distance from the singularities.
\end{corollary}

\subsection{The connection formula at one}

For nonresonant parameters, define
\begin{align*}
 A&=\frac{\Gamma(c)\Gamma(c-a-b)}
         {\Gamma(c-a)\Gamma(c-b)},\\
 B&=\frac{\Gamma(c)\Gamma(a+b-c)}
         {\Gamma(a)\Gamma(b)}.
\end{align*}
The classical connection formula is
\begin{equation}\label{eq:connection-one}
\begin{split}
 \F a b c z={}&A\,
 \F a b{a+b+1-c}{1-z}\\
 &+B(1-z)^{c-a-b}
 \F{c-a}{c-b}{1+c-a-b}{1-z}.
\end{split}
\end{equation}

\begin{theorem}[Certified connection formula]\label{thm:connection}
Equation \eqref{eq:connection-one} is a \StatusC{} iPRA theorem schema under:
\begin{enumerate}
\item apartness from the gamma poles and from the resonant hyperplanes
\(c-a-b\in\Z\);
\item a branch of \((1-z)^{c-a-b}\);
\item a finite parameter-continuation chain from a chamber in which the
connection coefficients are determined by effective Gauss evaluations;
\item a path-continuation certificate in the \(z\)-plane.
\end{enumerate}
\end{theorem}

\begin{proof}
Both terms on the right solve the original differential equation by
\cref{prop:one-solutions}.  Thus it suffices to determine the two constants in
one nonresonant parameter chamber.  One coefficient is the regular limit at
\(z=1\) and equals \(A\) by \cref{thm:gauss}.  Applying Euler's
transformation exchanges the two exponents at one; in a chamber where the
transformed regular limit exists, Gauss's theorem gives \(B\).  To place both
calculations in one formal proof, choose a finite path in parameter space that
avoids poles and resonance, and use contiguous relations whenever an integer
shift is more convenient.  Gamma recurrence and the ODE uniqueness theorem
show that the constants continue along the same chain.  Finally continue in
\(z\) with \cref{thm:path-continuation}.

Every step has finite data.  What is conditional is the supplied path and its
margins; no unrestricted theorem that the complement of all singular
hyperplanes is path connected is used inside iPRA.
\end{proof}

\begin{remark}[Resonant connection formulas]
When \(c-a-b\in\Z\), gamma coefficients in
\eqref{eq:connection-one} have poles and the two Frobenius solutions coalesce.
Taking an effective parameter limit produces a logarithmic formula involving
finite harmonic sums and gamma derivatives.  This can be formalized for any
fixed resonance and truncation order, but a uniform library of all resonant
cases is not included.
\end{remark}

\subsection{Barnes's contour integral}

The Barnes integrand is
\[
 G(s)=
 \frac{\Gamma(a+s)\Gamma(b+s)\Gamma(-s)}{\Gamma(c+s)}
 (-z)^s.
\]
A vertical contour must separate the poles
\(s=0,1,2,\ldots\) from
\(s=-a-n,-b-n\).  This separation is not expressible by a picture alone.

\begin{contract}[Barnes contour data]\label{con:barnes}
Supply:
\begin{enumerate}
\item a rational vertical line \(\RePart s=\sigma\), or a finite polygonal
deformation, with positive distances from every relevant pole family;
\item a branch of \(\log(-z)\) with
\(|\ImPart\log(-z)|\le\pi-\eta\) for rational \(\eta>0\);
\item explicit vertical-strip gamma bounds of the form
\[
 \norm{G(\sigma+it)}\le C(1+|t|)^M e^{-\eta|t|};
\]
\item for a closing rectangle, an explicit bound on the far vertical side,
typically containing \(\norm z^N\) when \(\norm z<1\).
\end{enumerate}
The gamma/Stirling and Cauchy-residue packets in the repository provide the
intended lower-level ingredients.
\end{contract}

\begin{theorem}[Certified Barnes integral]\label{thm:barnes}
Under \cref{con:barnes},
\begin{equation}\label{eq:barnes}
 \frac1{2\pi i}\int_{-i\infty}^{i\infty}G(s)\dd s
 =\frac{\Gamma(a)\Gamma(b)}{\Gamma(c)}\F a b c z
\end{equation}
with the contour interpreted as the supplied separating path.  The equality
has a primitive recursive precision schedule.
\end{theorem}

\begin{proof}
First truncate the contour at height \(T\) and close it after the first
\(N\) poles of \(\Gamma(-s)\).  The finite residue theorem gives a finite sum
of residues.  At \(s=n\), the residue, including the orientation sign, is
\[
 \frac{\Gamma(a+n)\Gamma(b+n)}{\Gamma(c+n)n!}z^n
 =\frac{\Gamma(a)\Gamma(b)}{\Gamma(c)}
   \frac{\poch a n\poch b n}{\poch c n n!}z^n.
\]
This is an exact gamma-recurrence calculation.  The horizontal edges and the
infinite vertical tails are bounded by
\(C(1+T)^Me^{-\eta T}\); bounded search chooses \(T\).  The far vertical edge
and the uncollected series tail are bounded by the supplied \(N\)-estimate;
choose \(N\) next.  The total error is allocated among the finitely many
edges and the series tail.  Thus the ordinary contour proof becomes a finite
residue identity plus primitive recursive inequalities.
\end{proof}

\begin{audit}
The report specifies the complete proof architecture but does not contain a
kernel-checked implementation of all complex gamma strip bounds and contour
integration.  Accordingly Barnes's formula is counted as \StatusC{}, not as
an unconditional finished analytic theorem.
\end{audit}

\subsection{Local monodromy and Schwarz ratios}

For a nonresonant Frobenius basis at a singular point with exponents
\(\rho_1,\rho_2\), continuation once around a small certified polygonal loop
multiplies the basis by
\[
 \begin{pmatrix}e^{2\pi i\rho_1}&0\\0&e^{2\pi i\rho_2}\end{pmatrix}.
\]
The statement is effective for rational or algebraically named exponents,
because the exponential has a Taylor modulus and the loop is finite.

\begin{proposition}[Möbius action on a solution ratio]\label{prop:schwarz-local}
Let \(s=\varphi_1/\varphi_0\) on a disk where \(\varphi_0\) is apart from
zero.  If continuation changes the basis by
\(\binom{\varphi_0}{\varphi_1}\mapsto
 M\binom{\varphi_0}{\varphi_1}\), then
\[
 s\longmapsto\frac{M_{10}+M_{11}s}{M_{00}+M_{01}s}.
\]
Moreover,
\[
 s'=\frac{W(\varphi_0,\varphi_1)}{\varphi_0^2},
\]
so local conformality is certified by apartness of the Wronskian and
\(\varphi_0\).
\end{proposition}

\begin{proof}
Both identities are one-line cross-multiplications.  The Wronskian has the
explicit form of \cref{lem:wronskian}, so its nonvanishing on a branch disk is
quantitative once the leading constant is apart from zero.
\end{proof}

\begin{audit}[Global Schwarz and monodromy claims]
The source page describes triangle maps, explicit monodromy matrices, and the
classification saying that certain rational exponent differences give finite
monodromy exactly when \(1/k+1/l+1/m>1\).  The local matrices and their finite
products are arithmetic objects and can be checked once connection matrices
are certified.  The global conformal mapping theorem and the complete finite
group classification are not reconstructed here.  They remain \StatusD{}.
\end{audit}

\subsection{The modular-lambda statement}

From \cref{cor:elliptic},
\[
 \F{1/2}{1/2}1z=\frac2\pi K(\sqrt z).
\]
Hence the ratio printed on the source page is
\[
 \tau=i\frac{K(\sqrt{1-z})}{K(\sqrt z)}
\]
for compatible square-root branches.  The assertion
\(\lambda(\tau)=z\), where lambda is a quotient of theta series, is the
classical elliptic-modular inversion theorem.

\begin{audit}[Modular inversion]
Theta series at an explicitly named \(\tau\) have elementary Gaussian tail
moduli, so their values are representable in iPRA.  What is missing is a
complete constructive proof that the theta quotient inverts the elliptic
integral ratio for every certified domain point, including the modular
transformation and global branch bookkeeping.  The package therefore proves
the reduction to the elliptic-modular theorem but marks the inversion itself
\StatusD{}.  The same missing modular infrastructure explains the status of
several isolated MathWorld values in \cref{audit:isolated-values}.
\end{audit}

\subsection{The John-transform sentence}

The Wikipedia page states that the Gauss function can be written as a John
transform but does not display a formula or normalization.  A mathematical
statement cannot be formalized from a bare reference to a transform with
multiple conventions.  Once an explicit kernel, domain, and normalization
are supplied, the intended proof is expected to be a finite change of
variables reducing the transform to Euler's integral.  No such theorem is
counted in the coverage total.
